\documentclass[11pt]{article}

\usepackage{amsmath,amssymb}
\usepackage{xurl}
\usepackage[colorlinks=true,linkcolor=blue,citecolor=blue,urlcolor=blue]{hyperref}

\input{./command.tex}

\title{Complete-Positivity Scope of the Peripheral Cyclicity Theorem}
\author{}
\date{2026-08-09}

\begin{document}
\maketitle
\gapnote{scope-restriction}{open}

\section{The source theorem}

Let $D\geq 1$, and let $T:\MN{D}\to\MN{D}$ be an irreducible positive unital
linear map satisfying the Schwarz inequality. Theorem~6.6 of
Ref.~\cite{Wolf2012Quantum}, stated at lines~244--252 of
\path{Notes/WolfNotePDF/wolf_ch6_sections.txt}, describes the peripheral
spectrum of $T$. Item~(1) states that this spectrum is a finite cyclic group
whose order $m$ satisfies
\begin{align}
    m
        &\leq D^2.
    \label{eq:wolf_thm6_6_kraus_scope_order_bound}
\end{align}
Items~(2)--(4) establish nondegeneracy, the existence of a peripheral unitary,
and a family of cyclic spectral projections, respectively
\cite[Theorem~6.6]{Wolf2012Quantum}. Neither complete positivity nor a Kraus
representation is a hypothesis of the source theorem.

\section{The formal restriction}

Let $I$ be a finite index set, let $K_i\in\MN{D}$ for each $i\in I$, and
define the Kraus map
\begin{align}
    E(X)
        &= \sum_{i\in I} K_i X K_i^\dagger.
    \label{eq:wolf_thm6_6_kraus_scope_kraus_map}
\end{align}
The available formal theorem is stated for this finite matrix family. It
assumes that $E$ is unital and irreducible and that the adjoint $E^\dagger$
has a positive-definite fixed point.\footnote{Formal declaration:
\leanid{Kraus.peripheralEigenvalues_eq_range_primitiveRoot} in
\path{QICLean/Channel/Peripheral/CyclicGroupKraus.lean}.}

The finite Kraus representation in
(\ref{eq:wolf_thm6_6_kraus_scope_kraus_map}) makes $E$ completely positive.
Complete positivity is stronger than the positive Schwarz hypothesis in
Theorem~6.6 of Ref.~\cite{Wolf2012Quantum}. The positive-definite fixed point
of $E^\dagger$ is also additional signature-level data; the source theorem
does not assume such a fixed point. In the finite-dimensional Kraus setting,
this datum should instead be derived in two steps. Unitality of $E$ makes
$E^\dagger$ trace preserving, which yields a positive-semidefinite fixed
point, and irreducibility then makes that fixed point positive definite.

The formal theorem proves the completely positive specialization of the
cyclicity conclusion in Theorem~6.6(1) of Ref.~\cite{Wolf2012Quantum}. It does
not prove the order bound in
(\ref{eq:wolf_thm6_6_kraus_scope_order_bound}), and it does not state
items~(2)--(4). Separate finite-Kraus companion declarations establish
nondegeneracy, peripheral unitaries, and cyclic projections under corresponding
Kraus hypotheses. These declarations do not provide the abstract positive
unital Schwarz-map formulation of the source theorem. The corresponding
Blueprint entry therefore uses a specialization title.

\section{Elimination plan}

The first step is to remove the explicit fixed-point assumption from the
finite-Kraus theorem. One must derive a positive-semidefinite fixed point of
$E^\dagger$ from trace preservation and then use irreducibility to prove that
the fixed point is positive definite.

The second step is to formulate peripheral eigenvectors and their
multiplicative-domain relations for an abstract positive unital Schwarz map.
The proof must establish the invertibility of peripheral eigenvectors and the
closure of peripheral eigenvalues under multiplication without selecting
Kraus operators. The finite-group argument used by the present cyclicity
theorem then applies without change. After this abstract result is available,
the finite-Kraus theorem should remain as a specialization rather than serve as
the source-facing form of Theorem~6.6 of Ref.~\cite{Wolf2012Quantum}.

\section{Verdict}

\textbf{Verdict: scope restriction.} The cyclicity declaration proves
Theorem~6.6(1) of Ref.~\cite{Wolf2012Quantum} under the stronger
complete-positivity hypothesis supplied by a finite Kraus representation. It
also takes a positive-definite adjoint fixed point as explicit input and omits
the bound in (\ref{eq:wolf_thm6_6_kraus_scope_order_bound}). Finite-Kraus
companion specializations cover conclusions corresponding to items~(2)--(4),
but the abstract Schwarz-map formulation of all four items and the order bound
in item~(1) remain open at the source-facing interface.

\bibliographystyle{alpha}
\bibliography{references}

\end{document}
